%============================================================
% HALAMAN KATA PENGANTAR
% Gunakan bahasa formal. Sertakan ucapan terima kasih kepada
% pihak-pihak yang berkontribusi pada kerja praktik.
%============================================================
\cleardoublepage
\phantomsection
\addcontentsline{toc}{chapter}{KATA PENGANTAR}
\begin{center}
  {\fontsize{12}{15}\bfseries KATA PENGANTAR\par}
\end{center}
\vspace{12pt}

Puji Syukur kami ucapkan kepada Tuhan Yang Maha Esa atas limpahan berkat dan rahmat-Nya
sehingga kami dapat menyelesaikan Laporan Kerja Praktik di \NamaPerusahaan\ yang
dilaksanakan mulai tanggal \PeriodeKP\ dengan judul \textit{``\JudulKP''}.

Laporan ini merupakan salah satu syarat bagi penyelesaian kerja praktik jurusan Teknik
Informatika PENS (Politeknik Elektronika Negeri Surabaya), dan juga sebagai sarana
pembelajaran dan pendalaman untuk menyesuaikan ilmu pengetahuan teoritis yang diterima di
perkuliahan dengan pengetahuan nyata di dunia kerja.

Selama proses pembuatan laporan kerja praktik, banyak bantuan, saran, dan motivasi yang
diberikan oleh berbagai pihak sehingga kami dapat menyelesaikan laporan ini. Oleh karena
itu, kami mengucapkan terima kasih kepada:

\begin{enumerate}[label=\arabic*., leftmargin=*, itemsep=0pt]
  \item Allah SWT, Tuhan Yang Maha Esa yang telah melimpahkan kemudahan kepada penulis
        sejak dalam pelaksanaan tugas kerja praktik di \NamaPerusahaan\ sampai dengan
        penyusunan laporan kerja praktik.
  \item Orangtua dan keluarga yang selalu memberikan dukungan moril maupun materiil sejak
        kerja praktik hingga penyusunan laporan kerja praktik.
  \item Bapak M. Udin Harun Al Rasyid, S.Kom., Ph.D. selaku Kepala Departemen Teknik
        Informatika dan Komputer PENS.
  \item Bapak \NamaKetuaProdi\ selaku Ketua Prodi \JenjangProdi\ \NamaProdi\ PENS.
  \item Ibu \NamaKoordinatorKP\ selaku Koordinator Kerja Praktik Jurusan Teknik Informatika PENS.
  \item Ibu \NamaPembimbing\ selaku Dosen Pembimbing yang telah meluangkan waktu dan
        memberi bimbingan, saran, dan arahan dalam proses pengerjaan proyek maupun saat KP
        berlangsung.
  \item \NamaPerusahaan\ yang telah memberikan kesempatan pada penulis untuk melaksanakan
        kerja praktik.
  \item Bapak \NamaPembimbingLapangan\ selaku \JabatanPembimbingLapangan\ sekaligus
        Pembimbing Lapangan di \NamaPerusahaan.
  \item Para karyawan/karyawati di \NamaPerusahaan\ yang telah bersedia bekerja sama dan
        memberikan informasi-informasi yang dibutuhkan dalam pelaksanaan kerja praktik ini.
  \item Serta pihak-pihak lain yang telah membantu kami namun tidak dapat disebutkan satu
        persatu.
\end{enumerate}

Demikian laporan kerja praktik ini disusun. Kami menyadari masih banyak kekurangan dalam
penyusunan laporan ini. Kami berharap semoga apa yang kami kerjakan dan kami tulis dalam
laporan ini dapat membantu memberi manfaat bagi pihak-pihak terkait.

\vspace{1cm}
\begin{flushright}
  \TanggalPengesahan\\[0.8cm]
  Penulis
\end{flushright}
